\documentclass[11pt]{article}
\usepackage[margin=0.78in]{geometry}
\usepackage{amsmath,amssymb,bm}
\usepackage{booktabs}
\usepackage{enumitem}
\usepackage{hyperref}
\hypersetup{colorlinks=true, linkcolor=black, urlcolor=black}
\setlength{\parindent}{0pt}
\setlength{\parskip}{0.55em}

\title{Cyrus Guided Notes: Stanford CS231A Computer Vision, From 3D Reconstruction to Recognition}
\author{Stanford course pack}
\date{2026-05-15}

\begin{document}
\maketitle

\section*{Course source}
\textbf{Official source:} \url{https://web.stanford.edu/class/cs231a/}

\textbf{Why this course is in the system.} Entry point for camera models, projection, epipolar geometry, optical flow, estimation, and 3D reconstruction.

\textbf{Learning output.} Already connected to local course notes and a guided PDF.

\section*{Prerequisite checkpoint}
\begin{itemize}[leftmargin=1.4em]
  \item Linear algebra
  \item Projective geometry
  \item Probability and estimation
\end{itemize}

\section*{Formula checkpoint}
Do not memorize these first. Read each formula as a sentence: what changes, what is observed, what is optimized, and what output it produces.
\begin{align*}
  s\begin{bmatrix}u\\v\\1\end{bmatrix}=K[R\;\bm{t}]\begin{bmatrix}X\\Y\\Z\\1\end{bmatrix}\\
  \bm{x}'^{\mathsf T}F\bm{x}=0\\
  E=[\bm{t}]_{\times}R\\
  \min_{T_i,P_j}\sum_{i,j}\rho(\|\bm{u}_{ij}-\pi(T_iP_j)\|_2^2)
\end{align*}

\section*{Visual page}
Show a 3D point, camera pose, image plane, and projection ray.

\section*{GoodNotes pages}
\begin{itemize}[leftmargin=1.4em]
  \item Page ST-CV001: camera projection
  \item Page ST-CV002: epipolar geometry
  \item Page ST-CV003: bundle adjustment
\end{itemize}

\section*{Web interaction}
A draggable projection lab showing how depth changes image coordinates.

\section*{Obsidian and Notion}
\begin{tabular}{p{0.22\linewidth}p{0.68\linewidth}}
\toprule
Layer & Output \\
\midrule
Obsidian & Stanford Spatial AI -> CS231A \\
Notion & Geometry notes and PDF evidence. \\
\bottomrule
\end{tabular}

\section*{First study move}
Open the official source, copy the prerequisite checkpoint into GoodNotes, then write one formula in your own words before watching or reading more.

\end{document}
